% paper/sections/appendix_stability_analysis.tex
%
% TVD-RK3 + Dissipative CCD のフォン・ノイマン安定性解析（§5 本文から移動）

\subsection{TVD-RK3 + Dissipative CCD の安定性解析}
\label{app:stability_derivation}

本節では，§~\ref{sec:advection_dccd_design} の Result~\ref{result:dccd_stability} で
引用した CFL 上界 $\sigma_{\max} \approx 0.91$ の導出を示す．

\medskip
\noindent\textbf{増幅係数の導出．}
TVD-RK3（式~\eqref{eq:tvd_rk3}）の増幅係数は，
線形空間演算子の固有値 $\lambda_j = -\mathrm{i}\,u\,\hat{k}^*_j H_j/h$（純虚数；
$\hat{k}^*_j = \hat{k}^*(\xi_j)$ は無次元 CCD 修正波数，$H_j \equiv H(\xi_j;\varepsilon_d^{\mathrm{adv}})$）
に対して
\[
  g_{\mathrm{RK3}}(-\mathrm{i}\omega) = 1-\mathrm{i}\omega - \frac{\omega^2}{2} + \frac{\mathrm{i}\omega^3}{6}
\]
（$\omega \equiv \sigma H(\xi)\hat{k}^*(\xi)$，$\sigma=u\Delta t/h$）であり，
$|g_{\mathrm{RK3}}|^2 = 1 - \omega^4/12 + \omega^6/36$
から安定条件 $|g| \leq 1$ は
\begin{equation}
  \omega^4\!\left(\frac{\omega^2}{36} - \frac{1}{12}\right) \leq 0
  \;\Longleftrightarrow\;
  |\omega| \leq \sqrt{3}
  \label{eq:rk3_stability_condition}
\end{equation}

\medskip
\noindent\textbf{CFL 上界の評価．}
式~\eqref{eq:rk3_stability_condition} から
$\sigma \leq \sqrt{3}/\max_\xi\bigl\{H(\xi;\varepsilon_d^{\mathrm{adv}})\,\hat{k}^*(\xi)\bigr\}$ となる．
無次元修正波数 $\hat{k}^*(\xi)$ は特性方程式（式~\eqref{eq:CCD_I} の Fourier 解析）から求まり，
本実装 ch13 既定の dissipative フィルタ強度 $\varepsilon_d^{\mathrm{adv}} = 0.05$
（CCD 移流再構成の散逸係数；式~\eqref{eq:dccd_adv_filter} で定義）のとき
$\max_\xi\{H(\xi)\,\hat{k}^*(\xi)\} \approx 1.9$
（$\xi \approx 0.8\pi$ 付近，再現スクリプト：
\texttt{scripts/dispersion\_dccd.py}；分散曲線図は別 CHK で添付予定）であるから：
\[
  \sigma_{\max} = \frac{\sqrt{3}}{1.9} \approx 0.91
\]

\medskip
\noindent\textbf{フィルタなしの場合との比較．}
純 CCD フィルタなし（$H=1$）の場合，
順 Euler 時間積分では $|g_\text{FE}|^2 = 1+\omega^2 > 1$（全波数で不安定；式~\eqref{eq:ccd_adv_instability}）だが，
TVD-RK3 では $\sigma_{\max} \approx \sqrt{3}/\hat{k}^*_{\max} \approx \sqrt{3}/2.35 \approx 0.74$（安定）．
Dissipative フィルタは安定限界の小幅な緩和よりも，Gibbs 振動の抑制（高波数減衰）を主目的とする．
